\documentclass[12pt]{article}
\usepackage[utf8]{inputenc}
\usepackage[T1]{fontenc}
\usepackage{hyperref}
\title{Yong Language Resources for K-12 Curriculum Development}
\author{Compiled by Claude Sonnet 5.5 (claude-sonnet-5-5) for Andr\'as Kornai}
\date{}
\begin{document}
\maketitle

\section{Introduction}
Yong (Tai Yong, ISO 639-3 yno) is a Southwestern Tai language of northern Thailand. Speaker estimates vary: 12,600 (2000) in one source and 14,000 in another. The Tai Yong trace their ancestry to Tai Lue people from Xishuangbanna and Kengtung, and the largest resettlement came in 1805 from Muang Yong. Communities are in Lamphun (Muang, Pasang, Ban Hong and Mae Tha districts), Lampang, Chiang Mai and Chiang Rai. The searches did not name Phayao as a location. Most speakers are also bilingual in Northern Thai and Standard Thai, and the language is reported as phonologically similar to Tai Lue. It is not known to be taught in schools. Very little educational material was found.

\section{Primary Educational Resources}
None were found in the search results.

\section{Linguistic and Reference Materials}
\begin{itemize}
\item Wikipedia on the Yong language: \url{https://en.wikipedia.org/wiki/Yong_language}
\item Wikipedia on the Kra--Dai languages: \url{https://en.wikipedia.org/wiki/Kra%E2%80%93Dai_languages}
\item A Thai university repository document (Srinakharinwirot University) appeared in the results. Its content was not verified: \url{https://ir.swu.ac.th/bitstreams/64bd6a5e-fba7-4cba-b430-88336c719564/download}
\end{itemize}

\section{Religious and Community Resources}
Joshua Project pages on the Yong people and language, which are missionary-oriented demographic profiles: \url{https://joshuaproject.net/people_groups/16053} and \url{https://joshuaproject.net/languages/yno}. A people-group profile PDF from Asia Harvest: \url{https://people-groups.asiaharvest.org/Thailand/Yong.pdf}

\section{Digital and Audio Resources}
None were found in the search results.

\section{Curriculum Development Context}
Yong is used by adults in the community but not by all young people, so intergenerational transmission is weakening. No orthography, school programme or published teaching material was found. Because of the closeness to Tai Lue and Northern Thai, the Tai script traditions of the north may be relevant, but this was not researched here. Community consultation would be the first step.
\end{document}